\label{chap:introduction}

\section{Problem Statement}
\label{sec_problem}

Firearms correspond to one of the most frequently used instruments worldwide for committing homicides. In forensic fieldwork, encountering gunshot wounds during external body examinations is common, whether at crime scenes or during autopsies. During these examinations, forensic experts must precisely determine the nature of the wound. The primary challenges involve distinguishing between entry and exit wounds and classifying the entry wounds according to the distance from which the shot was fired, an assessment commonly referred to as the Medical-Legal Shooting Distance (MLSD).

Under ideal conditions, certain morphological characteristics---such as shape, size, presence of an abrasion ring, and margins---allow for a relatively easy differentiation between entry and exit wounds. However, in reality, these characteristics are not always clearly visible; they may be partially absent, overlapping, or drastically altered by putrefaction. Other confounding factors, including the firearm caliber, ammunition type, shot angle, intermediate targets, and the impacted body region, further complicate accurate classification, even for experienced forensic pathologists.

The application of Artificial Intelligence (AI) and Deep Learning (DL) in computer vision tasks has recently revolutionized medical diagnostics, offering robust and automated tools for complex image classification. Despite these remarkable advancements in clinical settings, the translation of such technologies into the forensic domain remains extremely limited. One of the principal challenges lies in acquiring comprehensive datasets with high-quality, scientifically sound data. Forensic images are usually taken under various environmental conditions, with a range of photographic equipment, and within the chaotic, routine circumstances of crime scenes, making standardization practically impossible and creating highly imbalanced datasets.

Existing literature in the forensic domain, such as recent classification models trained on human autopsy databases, relies heavily on traditional Convolutional Neural Networks (CNNs). These architectures universally employ conventional Multi-Layer Perceptrons (MLPs) comprising dense linear layers for their final classification stages. While these architectures perform adequately under generalized circumstances, they inherently struggle to efficiently map the highly non-linear, complex feature boundaries present in noisy, imbalanced, and unstandardized forensic wound images. There is a critical gap and pressing need for a more advanced, efficient classifier that can extract and separate abstract non-linear patterns without requiring exponentially larger datasets or massive parameter counts.

\section{Research Objectives}
\label{sec_goals}

The primary objective of this study is to develop, implement, and evaluate a novel hybrid deep learning architecture, named \textbf{TransferKAN}, for the automated classification of human gunshot wounds using real crime scene images. This work focuses meticulously on two specific forensic tasks: classifying wound types (Entry vs. Exit) and determining the Medical-Legal Shooting Distance (MLSD) as Contact, Close range, or Distant.

To achieve this overarching goal, the following specific objectives were established:
\begin{itemize}
    \item To implement a Transfer Learning framework utilizing a pre-trained ResNet152 backbone, capitalizing on its proven capacity for high-level spatial feature extraction from complex images.
    \item To design and integrate a Kolmogorov-Arnold Network (KAN) as the primary decision classifier, replacing traditional dense MLP layers, in order to better capture and parameterize complex non-linear relationships through learnable splines.
    \item To implement a rigorous, case-based k-fold cross-validation scheme to evaluate the model's performance reliably, eliminating any possibility of semantic data leakage that typically inflates metrics in forensic AI research.
    \item To rigorously compare the performance metrics (Accuracy, Precision, Recall, F1-Score, AUC, and Specificity) of the proposed TransferKAN architecture against the baseline performance of state-of-the-art CNNs, specifically the ResNet152 baseline model reported in the 2024 literature.
\end{itemize}

\subsection{Research Hypothesis}
\label{sec_hypothesis}

We hypothesize that replacing the standard Multi-Layer Perceptron (MLP) classification head of a state-of-the-art deep convolutional network (ResNet152) with a Kolmogorov-Arnold Network (KAN) will result in a superior, highly robust hybrid architecture (TransferKAN). By learning activation functions on the network's edges rather than relying on fixed node activations, the TransferKAN architecture will better model the non-linear boundaries of imbalanced forensic image classes. Consequently, this approach will yield higher discriminative performance and reliability---particularly evidenced by the Area Under the Curve (AUC)---in both wound type and MLSD classification compared to the traditional ResNet152 baseline.

\section{Contributions}
\label{sec_contributions}

The main contributions of this work can be summarized as follows:
\begin{itemize}
    \item The pioneering introduction of the \textbf{TransferKAN architecture} into the forensic pathology domain. This hybrid model successfully leverages the deep feature-extraction prowess of ResNet152 alongside the parameter-efficient, dynamic non-linear classification capabilities of Kolmogorov-Arnold Networks.
    \item The empirical demonstration that the TransferKAN approach, coupled with robust data augmentation and strict k-fold cross-validation, decisively outperforms the traditional CNN baselines utilized in prior forensic literature, particularly showing vast improvements in robustness and minority class detection for Entry/Exit and MLSD classifications on a real-world crime scene dataset.
    \item The provision of a rigorous methodological pipeline that actively mitigates data leakage issues through case-based data splitting, offering a much more realistic evaluation of the model's true generalization capabilities in authentic forensic investigative scenarios.
\end{itemize}

\section{Organization of This Document}
\label{sec_about}

The remainder of this dissertation is organized as follows. Chapter~\ref{chap:background} provides the essential background on the forensic pathology of gunshot wounds and introduces the theoretical underpinnings of Deep Learning and Kolmogorov-Arnold Networks (KAN). Chapter~\ref{chap:related} critically reviews the related literature on AI applied to forensic wound classification, identifying current gaps. Chapter~\ref{chap:solution} introduces the proposed solution, detailing the TransferKAN architecture. Chapter~\ref{chap:methodology} describes the proposed methodology, including the dataset, data augmentation, and cross-validation procedures. Chapter~\ref{chap:results} presents a comprehensive comparative analysis of the results between TransferKAN and the established baselines. Finally, Chapter~\ref{chap:conclusions} concludes the thesis, summarizing the findings and suggesting potential avenues for future research.

